% Extended results discussion for Experiment 03: Catastrophic Model Failure.
% This file provides a standalone narrative of the catastrophic failure
% experiment results, suitable for inclusion in an appendix or supplementary
% material if a more detailed treatment is needed.

\subsection{Experiment~3: Catastrophic Failure --- Extended Discussion}
\label{appendix:catastrophic_failure_extended}

The catastrophic failure experiment evaluates ParetoBandit's resilience
when a previously top-performing model (Mistral-Large) begins returning
garbage responses mid-deployment.  This scenario---cloud API outages,
5xx errors, timeout cascades, and silent quality degradation after
provider-side model updates---is routine in production multi-LLM
deployments.

\paragraph{Three-phase design.}
The test split ($n{=}608$ prompts per phase) is divided into:
Phase~1 (normal operation), Phase~2 (Mistral failure: reward
$\to\cfFailureReward$, cost unchanged), and Phase~3 (Mistral recovered,
reusing Phase~1 prompts).  The model registry is \emph{not} mutated
at boundaries; adaptation is driven purely by observed rewards.

\paragraph{Key findings.}
Three algorithm variants are compared at each of three budget targets
(Tight, Moderate, Loose), plus an unconstrained baseline (20~seeds):

\begin{itemize}[nosep,leftmargin=*]
  \item \textbf{Naive Bandit ($\gamma{=}1.0$)} detects the failure
    and reduces Mistral traffic, recovering reward.  However, its
    static cost penalty cannot adjust to the changed cost landscape:
    at moderate budget, it under-spends at \cfNaiveModPhaseTwoRatio{}
    in Phase~2 as traffic shifts to cheaper Llama-8B, while at loose
    budget the Naive Bandit over-spends to fill the gap left by
    Mistral with expensive Gemini traffic.
  \item \textbf{Forgetting Bandit ($\gamma{=}\cfGamma$)} shares
    ParetoBandit's forgetting factor but replaces the closed-loop
    pacer with a static cost penalty.  It detects the failure faster
    than the Naive Bandit (geometric forgetting down-weights stale
    Phase~1 Mistral rewards), but budget compliance is \emph{worse}
    because forgetting also erases the cost penalty's prior anchoring
    effect.  At tight budget, the Forgetting Bandit overshoots by
    \cfForgetTightPhaseTwoRatio{} in Phase~2 and
    \cfForgetTightPhaseThreeRatio{} in Phase~3---confirming that
    \textbf{the BudgetPacer, not forgetting, is responsible for
    ParetoBandit's cost invariance}.
  \item \textbf{ParetoBandit ($\gamma{=}\cfGamma$)} detects the failure,
    redistributes traffic (Mistral drops to
    \cfParetoBanditModMistralPhaseTwo\% at moderate budget), \emph{and}
    maintains budget compliance: \cfParetoBanditModPhaseOneRatio,
    \cfParetoBanditModPhaseTwoRatio, and
    \cfParetoBanditModPhaseThreeRatio{} across the three phases.
\end{itemize}

\paragraph{Budget compliance is the key differentiator.}
ParetoBandit maintains budget ratios within $0.95$--$1.00\times$
across all twelve phase--budget combinations.  The Naive Bandit, despite
recovering reward, violates the budget by up to
\cfNaiveTightPhaseThreeRatio{} (tight, Phase~3).  The Forgetting Bandit
fares even worse: at tight budget, it reaches
\cfForgetTightPhaseThreeRatio{} in Phase~3---geometric forgetting erases
the cost penalty's anchoring effect, producing runaway spend toward
expensive models.  This demonstrates that online learning alone---even
with forgetting---is insufficient for budget-constrained non-stationary
routing: the dual-variable mechanism is essential for maintaining cost
invariance through disruptions.

\paragraph{Phase~2 cost mechanism.}
The pacer's closed-loop response to catastrophic failure involves a
\emph{silent quality regression}: Mistral's reward drops to
\cfFailureReward{} while its cost remains unchanged---the provider
continues charging full price for degraded responses.
The reward signal drives traffic away from Mistral toward Llama-8B
(cheap) and, at looser budgets, Gemini-Pro (expensive).  Because
Mistral's cost is intermediate (normalized $\tilde{c}{=}0.33$), the
net cost impact of redistribution depends on where traffic lands: shifts
toward Llama reduce cost, while shifts toward Gemini increase it.
The BudgetPacer continuously adjusts $\lambda_t$ to absorb these
routing-induced cost changes, maintaining compliance throughout the
failure window.
No other condition achieves this: the Naive and Forgetting Bandits
shift traffic without any mechanism to detect and correct the
resulting budget deviations.

\paragraph{Figure traces (Figure~\ref{fig:catastrophic_failure}).}
The three panels overlay four conditions: ParetoBandit at three
budget levels (solid coloured), the Naive Bandit and Forgetting
Bandit at moderate budget (dashed grey and orange), and the
unconstrained baseline (dash-dot green).  We discuss each
condition's expected behaviour across all panels.

\emph{Panel~(a): Gemini-Pro selection fraction.}
Under normal Phase~1 pricing, Gemini-Pro is the most expensive arm
($\tilde{c}{=}0.58$).  Budget-constrained conditions therefore show
near-zero Gemini adoption: the static penalty or dual variable
penalises Gemini's cost.  The unconstrained baseline, free of cost
pressure, selects Gemini when UCB exploration favours it.
When Mistral fails in Phase~2, all learning conditions redirect
traffic---but \emph{not} toward Gemini.  At moderate budget, the
static penalty ($\lambda_s{=}0.30$) still penalises Gemini
(whose pricing is unchanged), so the Naive Bandit and Forgetting
Bandit shift traffic primarily to Llama-8B (the only cheap,
functioning arm).  ParetoBandit also routes mostly to Llama at
tight budget, while at loose budget residual Gemini adoption
increases modestly because $\lambda_t$ is low.
The unconstrained baseline shifts ${\sim}80\%$ of traffic to
Gemini---the highest-quality surviving arm---driving the cost
spike visible in panel~(c).
In Phase~3, Gemini adoption returns to near Phase~1 levels for all
constrained conditions.  The unconstrained baseline retains elevated
Gemini selection because its routing is quality-dominated.

\emph{Panel~(b): Windowed mean reward.}
Phase~1 reward is stable across conditions (${\sim}0.85$--$0.93$).
At the Phase~2 boundary, all conditions suffer an initial reward
trough as residual Mistral traffic returns garbage (reward
$\to\cfFailureReward$).  The depth of this trough depends on how
much Mistral traffic remains:
the \textbf{unconstrained baseline} recovers fastest because it
freely shifts to high-quality Gemini;
\textbf{ParetoBandit} recovers within ${\sim}\cfParetoBanditModAdaptSteps$ steps
but must route to cheaper arms (primarily Llama), producing a lower
steady-state reward than unconstrained;
the \textbf{Naive Bandit} (grey dashed) recovers reward comparably
to ParetoBandit because it also diverts from Mistral, but at
unchecked cost;
the \textbf{Forgetting Bandit} (orange dashed) recovers faster
than the Naive Bandit because geometric forgetting discounts stale
Phase~1 Mistral rewards, enabling quicker re-routing.
In Phase~3, all adaptive conditions recover toward Phase~1 reward
levels.  The Forgetting Bandit recovers slightly faster than the
Naive Bandit because Phase~2 failure observations fade within the
${\sim}333$-step effective memory window ($1/(1{-}\gamma)$).

\emph{Panel~(c): Windowed cost.}
This panel exposes the critical difference.  ParetoBandit's cost
traces hug the budget target lines throughout all three phases:
the BudgetPacer continuously adjusts $\lambda_t$ to absorb the
Phase~2 cost shock (see Phase~2 cost mechanism above).
The Naive Bandit's cost (grey dashed) drops to
$\cfNaiveModPhaseTwoRatio$ of the moderate target in Phase~2:
traffic shifts from Mistral (intermediate cost) primarily to
Llama-8B (cheapest arm), and the static penalty cannot loosen
to exploit the freed budget for higher-quality routing.
In Phase~3, the Naive Bandit recovers to
$\cfNaiveModPhaseThreeRatio$ as infinite memory anchors the
routing mix, but this compliance is coincidental---the static
penalty offers no guarantee across budget regimes.
The Forgetting Bandit's cost (orange dashed) shows a similar
pattern: already overshooting in Phase~1
($\cfForgetModPhaseOneRatio$), spiking in Phase~2
($\cfForgetModPhaseTwoRatio$), and worsening in Phase~3
($\cfForgetModPhaseThreeRatio$).  Forgetting accelerates
failure detection but also erases the cost penalty's anchoring
effect, producing runaway spend toward expensive models.
The unconstrained baseline (dash-dot green) sits far above all
constrained conditions, confirming that it operates without any
cost regard.

\paragraph{Recovery dynamics.}
When Mistral is restored in Phase~3, geometric forgetting
($\gamma{=}\cfGamma$) allows re-discovery via UCB exploration
as Phase~2 failure memory is discounted.  Phase~3 reward
approaches but does not fully match Phase~1 levels within the
608-step phase---consistent with the ${\sim}333$-step effective
memory window requiring time to accumulate sufficient positive
evidence.

\paragraph{Phase~3 within-subject design.}
Phase~3 deliberately reuses the same prompts as Phase~1 so that the
P1-vs-P3 comparison is a controlled within-subject design: any
performance difference is attributable solely to the router's internal
state, not prompt-sampling variability.
A potential concern is that the bandit's sufficient statistics already
contain Phase~1 observations of these exact feature vectors, giving
Phase~3 a ``warm-cache'' advantage over evaluation on fresh prompts.
Three considerations mitigate this:
(i)~under $\gamma{=}\cfGamma$, each Phase~1 observation has decayed to
effective weight $\gamma^{1216} \approx 0.026$ by the start of Phase~3,
so 97\% of the Phase~1 information has been forgotten;
(ii)~LinUCB maintains a $d$-dimensional $\theta$ vector per arm learned
from \emph{all} observations---re-encountering the same PCA feature
vector provides no per-prompt advantage, as the model generalises over
the feature distribution rather than memorising individual contexts;
(iii)~all four conditions receive identical Phase~3 prompts, so any
residual warm-cache effect is a controlled confound that does not affect
relative comparisons between algorithms.
